%% kapa: exam=1 round=35 year=2024 period=1 form=A subject=economics subject_ko=경제학원론 seq=27 number=67
%% source: Q-Net 감정평가사 기출문제 (cst00309) artlSeq=5213229 file=raw/1008/5213229/1교시 문제지 원본.pdf
%% conversion: kordoc 4.24.51 -> tools/gen1.py (auto); review=proofread-sonnet; flags=none
\begin{question}{67}
아래와 같은 고전학파 모형에서 정부지출이 150에서 200으로 증가할 경우 실질이자율과 민간투자의 변화에 관한 설명으로 옳은 것은? (단, $S,\overline{Y},\overline{T},\overline{G},I,r,s(r)$은 각각 총저축, 총생산, 조세, 정부지출, 투자, 실질이자율(\%), 민간저축률이며, 민간저축률은 실질이자율의 함수이다.)
\begin{framedtext}
○ $S=s(r)(\overline{Y}-\overline{T})+(\overline{T}-\overline{G})$ \par
○ $I=200-10r$ \par
○ $\overline{Y}=1000,\ \overline{T}=200,\ \overline{G}=150$ \par
○ $s(r)=0.05r$ \par
\end{framedtext}
\begin{choices}
  \item 실질이자율은 1 \%포인트 상승하고 민간투자는 10 감소한다.
  \item 실질이자율은 3 \%포인트 상승하고 민간투자는 30 감소한다.
  \item 실질이자율은 5 \%포인트 상승하고 민간투자는 50 감소한다.
  \item 실질이자율과 민간투자는 변화가 없다.
  \item 실질이자율은 1 \%포인트 하락하고 민간투자는 10 증가한다.
\end{choices}
\end{question}
\answer{1}
